\section{Computing the exceptional integer and the complete prime function}
\label{sec:effective-exception}

The generic-flatness integer of B04 is not merely existential when the
arithmetic equations of the parameter space are available. The following
construction gives a computable integer, without claiming it is the
smallest possible one. It is an effective form of the use of
\cite[Theorem 6.9.1]{EGAIV2} in B04.

\begin{proposition}[R44: an explicit generic-flatness certificate]
\label{prop:R44}
Let $R=\mathcal O_K[1/b_0]$, where $K$ is an explicitly presented number
field and $b_0$ is a positive integer. Given a finite affine cover and
equations for a finite-type $R$-scheme $Y$, one can compute a positive
integer $b$ such that $Y_{R[1/b]}$ is flat over $R[1/b]$.

In particular, apply this algorithm to the bounded cycle parameter space
of B01, constructed effectively as in R42. The resulting $b$, enlarged
by $b_0$ and the chosen coefficient characteristic, is a valid
exceptional integer $b_D$ in B04. It depends on the arithmetic model,
polarization and degree bound $D$, and not on the target class or its
multiplier.
\end{proposition}

\begin{proof}
For each affine chart, take a presentation
\[
 R[x_1,\ldots,x_s]/I,\qquad I=(f_1,\ldots,f_a).
\]
Fix a monomial order. Compute a monic Gr\"obner basis
$G=\{g_1,\ldots,g_c\}$ of $I_K\subset K[x_1,\ldots,x_s]$,
recording polynomial expressions
\[
 g_i=\sum_j u_{ij}f_j,\qquad
 f_j=\sum_i v_{ji}g_i
\]
over $K$. A standard basis algorithm tracks the first expressions,
and polynomial division supplies the second. All lists are finite.
This is the usual effective polynomial-ideal computation; the
monic standard-basis formalism over coefficient rings is discussed
in \cite[\S3, Definitions 3.8--3.10 and Propositions 3.11--3.13,
pp.~15--17; proof of Theorem 3.6, pp.~17--18]{HaimanSturmfels}.

Use an integral basis of $\mathcal O_K$ to clear the coefficients of
the $g_i,u_{ij},v_{ji}$ by a positive rational integer $b_Y$.
Over $R'=R[1/b_Y]$ the two displayed identities remain valid and
$I_{R'}=(G)$. Division by the monic $G$ expresses every polynomial
over $R'$ as an $R'$-linear combination of standard monomials,
modulo $I_{R'}$. Those monomials are independent over $R'$:
any relation would become a relation between the standard monomials
over the fraction field $K$, where they form a basis. Since
$R'\hookrightarrow K$ is injective, every coefficient in the relation
is zero. Hence $R'[x]/I_{R'}$ is a free, possibly infinite-rank,
$R'$-module, and in particular is flat.

The generic empty-chart case is included: if $I_K=(1)$, the recorded
expression for $1$ makes the localized chart empty. If $I_K=(0)$,
the empty standard basis gives the polynomial algebra and no additional
denominator is needed.

Take the product of the finitely many integers $b_Y$ obtained from the
charts. Flatness is local on the source, so this product works for $Y$.
One may apply the construction either to the pair space
$\mathcal H_D$ or to its individual-cycle factor $\mathcal Q_D$;
in the latter case its self-product is flat as well, by base change
and composition of flat morphisms. B04 then proves the asserted
lifting property. The algorithm has made no reference to a class
condition or a positive integer multiplier.
\end{proof}

\begin{corollary}[R45: an effective finite output for the local minimum]
\label{cor:R45}
For an explicitly presented arithmetic Eisenstein Weil fourfold,
with an extending polarization, a specified good-reduction model,
and the exact Betti-class input of R43, the function
\[
                    v\longmapsto c_{H_v}(t_v)
\]
on the geometric fibers of that good model is computable as a generic
integer $C_0$ together with a finite table of exceptional values.
The procedure terminates under the cited theorem inputs and the
explicit family bound of R24. No bound for its running time is asserted.
\end{corollary}

\begin{proof}
R43 computes $C_0$ and an actual rational cycle representing the
generic target, with multiplier and numerator degrees at most $C_0$.
The computation uses the supplied exact Betti marking; it does not
guess a period from numerical approximations. Enlarge the number field
to define the cycle and its supports when necessary, and use the
corresponding base change of the given good model. The class then has
a specified compatible specialization at every place in question.
This finite extension changes neither the generic geometric minimum
nor the geometric minimum at a fixed chosen geometric place.

If $C_0=1$, all the good local minima equal one. Otherwise R42
constructs the relative bounded parameter space for $D=C_0-1$, and
R44 computes a valid exceptional integer. Factor it, together with
the finitely many excluded rational characteristics. There are only
finitely many places of the defining number field above these primes.
At every other place B04 and specialization give exact equality with
$C_0$.

At each remaining place belonging to the given good model, specialize
the already computed generic rational cycle. R42 constructs its
finite-field bounded parameter spaces. R40, with upper bound $C_0$,
returns the exact local minimum and cycles attaining it. This produces
the finite table, as in R41. If a larger number field was used to
define the representative, the output is attached to the corresponding
specified geometric places; it is not an assertion that a class lacking
descent has a well-defined specialization at an unmarked smaller-field
place.
\end{proof}

Thus the geometric and arithmetic gates are effective in principle
when their actual input data are supplied. For the two explicit product
classes of this paper the table is empty and $C_0=16$, by P13 and V06,
without running these very large general computations. T5 itself gives
no equations for a different arithmetic member or numerical coordinates
for its Weil vector, so no further numerical table is asserted for an
unspecified member. R45 is the precise algorithmic form of that more
general question.
